\begin{table}[t]
\centering
\small
\caption{Negative and diagnostic results. These rows are reported as limitations rather than removed.}
\label{tab:limitations}
\begin{tabularx}{\linewidth}{lllXrrr}
\toprule
Case & Cert. & Status & Quality & Gain (pct.) & Calib. JVP & Held-out JVP \\
\midrule
M5 & CVaR & accepted & accepted gain, risk exceeds target & 4.7 & 0.014 & 0.353 \\
M5 & Bernstein & fallback & certificate mismatch & 0.0 & 0.718 & 0.885 \\
Power & CVaR & fallback & no accepted gain & 0.0 & 0.010 & 0.070 \\
Power & Cantelli & failed & over-conservative & 0.0 & 0.130 & 0.073 \\
French & Cantelli & fallback & over-conservative & 0.0 & 0.000 & 0.000 \\
\bottomrule
\end{tabularx}
\vspace{2pt}
\begin{minipage}{.98\linewidth}
\footnotesize \textit{Note.} Quality identifies why the optimized row is not counted as primary positive evidence. Fallback rows report the accepted equal-allocation solution when the optimized candidate is infeasible, numerically unacceptable, or not economically better.
\end{minipage}
\end{table}
